% NOTE: Filename is fixed by the book outline; content teaches Week 21
% (Security, IAM & Networking) of the syllabus.
\chapter{Security, IAM, and Networking for Data Platforms}
\index{AWS IAM}\index{permissions boundary}\index{service control policy}\index{VPC endpoint}\index{AWS PrivateLink}\index{AWS KMS}\index{Secrets Manager}\index{data perimeter}%
\begin{objectives}
\begin{itemize}
    \item Apply least-privilege IAM: roles, policies, permissions boundaries, and service control policies (SCPs).
    \item Keep data traffic off the public internet with gateway and interface VPC endpoints and PrivateLink.
    \item Manage encryption keys with AWS KMS and application secrets with Secrets Manager, including rotation.
    \item Build and verify an S3 gateway endpoint and access a bucket from EC2 without an internet gateway.
    \item Design an anti-exfiltration data perimeter that holds even when IAM credentials are compromised.
    \item Reason about the cost of the perimeter itself: endpoint-hours, KMS requests, and budget alarms.
\end{itemize}
\end{objectives}

\begin{crosscloud}
Identity, private networking, and key management form the security substrate of every cloud; the concepts are identical, the acronyms differ.

\vspace{2mm}
{\footnotesize
\begin{tabular}{@{}p{2.8cm}p{3.0cm}p{3.0cm}p{3.0cm}@{}}
\toprule
\textbf{Capability} & \textbf{AWS} & \textbf{Azure} & \textbf{GCP} \\
\midrule
Identity and access & IAM roles/policies & Entra ID + Azure RBAC & Cloud IAM \\
Org-wide guardrails & SCPs (Organizations) & Azure Policy & Organization Policy \\
Permissions ceiling & Permissions boundaries & Management groups + deny & IAM deny policies \\
Private service access & VPC endpoints / PrivateLink & Private Endpoints / Private Link & Private Service Connect \\
Key management & AWS KMS & Azure Key Vault & Cloud KMS \\
Secrets management & Secrets Manager & Key Vault secrets & Secret Manager \\
\addlinespace
Free S3-class endpoint & Gateway endpoint & Service endpoints & Private Google Access \\
Anti-exfiltration primitive & Endpoint policy + SCP + aws:PrincipalOrgID & Private Link + policy & VPC Service Controls \\
\bottomrule
\end{tabular}}

\vspace{1mm}
\textbf{Snowflake} adds account-level network policies and external-private-endpoint rules; \textbf{MongoDB} Atlas uses private endpoints and IP access lists; \textbf{Fabric} centralizes on Entra conditional access and tenant-level perimeter settings. The pattern---identity plus network plus keys---is universal.
\end{crosscloud}

\section{Week 21 Roadmap and Mental Model}
Every chapter so far ended with an implicit question: who is allowed to touch this, and over what network path? Week~21 makes the answers explicit. The week's stack has three layers: \textbf{identity} (IAM roles and policies, bounded by permissions boundaries and SCPs), \textbf{network} (VPC endpoints and PrivateLink so traffic never leaves the AWS backbone), and \textbf{cryptography} (KMS keys and Secrets Manager so access requires more than a stolen credential).

The mental model that organizes all three is the \emph{data perimeter}: a set of controls ensuring that data can only be accessed by expected identities, from expected networks, within the expected organization \cite{awsIAMDocs}. A data platform that is fast, cheap, and exfiltratable is a liability; this week is how you make it merely fast and cheap.

\begin{definitionbox}
\textbf{A data perimeter} is the combination of identity, network, and organization-boundary controls that constrains data access to expected principals, expected networks, and the expected organization---designed to hold even if a credential is stolen.
\end{definitionbox}

\begin{keyterms}
\textbf{Role}: an assumable identity with temporary credentials---the default way workloads authenticate. \textbf{Permissions boundary}: a maximum-permissions ceiling on a role, regardless of its policies. \textbf{SCP}: an Organizations-level guardrail bounding what accounts may do. \textbf{Gateway endpoint}: free VPC route-table entry for S3/DynamoDB. \textbf{Interface endpoint}: PrivateLink-backed ENI for most other services, billed per AZ-hour plus data. \textbf{KMS}: managed key service with envelope encryption and key policies. \textbf{Secrets Manager}: secrets storage with automatic rotation.
\end{keyterms}

\section{Monday: IAM Deep Dive---Roles, Boundaries, and SCPs}
\subsection{Least Privilege as an Engineering Practice}
The unit of AWS authorization is the \emph{role}: workloads assume roles and receive temporary credentials---never long-lived access keys. Least privilege is an iterative practice, not a one-time design: start from the actions the workload provably needs (from CloudTrail or Access Analyzer), scope resources by ARN and prefix, add conditions (\texttt{aws:SourceVpc}, \texttt{aws:PrincipalTag}), and re-review on a cadence \cite{awsIAMDocs}.
\index{least privilege}\index{IAM role}

\subsection{Ceilings: Permissions Boundaries and SCPs}
Two mechanisms bound what even a misconfigured policy can grant. A \emph{permissions boundary} caps a role's maximum permissions---delegated role-creation (developers create their own Glue job roles) is safe only because the boundary prevents those roles from exceeding their lane. An \emph{SCP} applies at the AWS Organizations level and cannot be overridden by any account administrator: ``deny all S3 access from outside our VPC endpoints,'' ``deny leaving the organization,'' ``restrict regions.'' Boundaries protect role hygiene; SCPs protect the organization.
\index{permissions boundary}\index{service control policy}

\begin{pitfall}
\textbf{An allow without a ceiling.} IAM evaluates identity policies \emph{and} boundaries \emph{and} SCPs---the effective permission is the intersection. Teams that only write identity policies leave both ceiling mechanisms empty, so any future over-broad grant is instantly effective. Deploy the ceilings first; they are cheap and invisible until the day they save you.
\end{pitfall}

\section{Tuesday: VPC Endpoints and PrivateLink}
\subsection{Gateway Versus Interface}
A \emph{gateway endpoint} (S3 and DynamoDB only) adds a route-table entry directing service traffic over the AWS backbone---no ENI, no hourly charge, and an endpoint policy that can restrict which buckets are reachable. An \emph{interface endpoint} places an elastic network interface in your subnet backed by AWS PrivateLink, available for most AWS services and for third-party/private services---billed per AZ-hour plus per-GB processed \cite{awsPrivateLinkDocs}.
\index{gateway endpoint}\index{interface endpoint}\index{PrivateLink}

The data-platform consequence: with an S3 gateway endpoint, a VPC with \emph{no internet gateway at all} can still reach the lake---which is exactly the posture regulated workloads want. Glue, Redshift, Athena, KMS, and the rest get interface endpoints; traffic to AWS APIs never touches the public internet.

\begin{notebox}
\textbf{Cost guardrail:} interface endpoints bill per AZ-hour plus per-GB. For high-volume S3 traffic always prefer the free gateway endpoint, and consolidate interface endpoints (one shared-services VPC) instead of one per workload VPC. Set AWS Budgets alerts per account \emph{before} the hardening project, so the perimeter itself never produces a surprise bill.
\end{notebox}

\subsection{Endpoint Policies: The Network Half of the Perimeter}
An endpoint policy restricts what the endpoint will carry: ``this S3 endpoint serves only these buckets, only for principals in our organization.'' Endpoint policies are how the network layer enforces the data perimeter independently of IAM---a stolen credential used from a laptop cannot route through your endpoint, and your endpoint will not carry data to a stranger's bucket.
\index{endpoint policy}

\section{Wednesday: KMS and Secrets Manager}
\subsection{KMS: Keys as Policy Objects}
AWS KMS provides envelope encryption: data keys encrypt data; KMS customer master keys encrypt data keys \cite{awsKMSDocs}. For the data engineer, the essentials are: customer-managed keys (CMKs) when you need key policy control, rotation evidence, and cross-account sharing (the lake's key must allow the analytics account to decrypt); per-domain keys to bound blast radius; and CloudTrail logging of every key use, which makes ``who decrypted what'' an answerable question. Grants and key policies together decide whether a stolen credential can decrypt anything---with a per-service, per-bucket key policy, often it cannot.
\index{KMS}\index{envelope encryption}\index{customer-managed key}

\subsection{Secrets Manager: No More Static Credentials}
Database passwords, API keys, and JDBC strings live in AWS Secrets Manager, retrieved at runtime by IAM-authenticated callers, and---the point---\emph{rotated automatically} on a schedule with Lambda rotators that update both the secret and the target database \cite{awsSecretsManagerDocs}. Static credentials in environment variables, DAG code, or \texttt{.env} files are the most common real-world data-platform breach vector; rotation converts a leaked credential from a breach into an incident ticket.
\index{Secrets Manager}\index{rotation}

\begin{tipbox}
Redshift, RDS, and document databases have native Secrets Manager rotation integrations. For anything else, the pattern is: secret in Secrets Manager, application fetches at startup and caches briefly, rotation Lambda updates both sides, and CloudTrail alarms on unusual \texttt{GetSecretValue} patterns.
\end{tipbox}

\section{Thursday Hands-on Project: Private S3 Access with No Internet Route}
\index{hands-on project!VPC endpoint for S3}
Create a VPC endpoint for S3, then access a bucket from an EC2 instance in a VPC with no internet gateway---proving the private path works and the public path does not exist.

\subsection*{Lab Objectives}
\begin{itemize}
    \item Build a VPC with a private subnet and no internet gateway.
    \item Add an S3 gateway endpoint with a restrictive endpoint policy.
    \item From EC2 (via Systems Manager Session Manager, since there is no inbound path), copy a file to the bucket.
    \item Demonstrate that the endpoint policy blocks access to a non-authorized bucket.
\end{itemize}

\subsection*{Procedure}
\begin{lstlisting}[language=bash,caption={Private VPC, gateway endpoint, and the proof.},label={lst:ch21-endpoint}]
# 1. VPC with a private subnet; deliberately no IGW is created or attached.
$VpcId = (aws ec2 create-vpc --cidr-block 10.21.0.0/16 `
  --query Vpc.VpcId --output text)
$SubnetId = (aws ec2 create-subnet --vpc-id $VpcId `
  --cidr-block 10.21.1.0/24 --query Subnet.SubnetId --output text)
$RtId = (aws ec2 describe-route-tables --filters Name=vpc-id,Values=$VpcId `
  --query 'RouteTables[0].RouteTableId' --output text)

# 2. Gateway endpoint for S3, scoped by endpoint policy to one bucket.
$Policy = @"
{
  "Version": "2012-10-17",
  "Statement": [{
    "Effect": "Allow",
    "Principal": "*",
    "Action": ["s3:GetObject", "s3:PutObject", "s3:ListBucket"],
    "Resource": [
      "arn:aws:s3:::week21-private-bucket-123456789012",
      "arn:aws:s3:::week21-private-bucket-123456789012/*"
    ]
  }]
}
"@
$Policy | Set-Content -Path endpoint-policy.json -Encoding ascii

aws ec2 create-vpc-endpoint --vpc-id $VpcId `
  --service-name com.amazonaws.us-east-1.s3 `
  --route-table-ids $RtId `
  --policy-document file://endpoint-policy.json

# 3. EC2 in the private subnet (SSM-enabled role; no public IP, no IGW).
aws ec2 run-instances --image-id ami-0abcdef1234567890 `
  --instance-type t3.micro --subnet-id $SubnetId `
  --iam-instance-profile Name=Week21SSMRole --no-associate-public-ip-address

# 4. From the instance (via Session Manager):
#    aws s3 cp test.txt s3://week21-private-bucket-123456789012/   # succeeds
#    aws s3 ls s3://some-other-bucket/                             # denied by endpoint policy
\end{lstlisting}

\subsection*{Verification}
\begin{itemize}
    \item The copy to the authorized bucket succeeds with no internet gateway attached---traffic used the gateway endpoint.
    \item Access to any other bucket fails with an explicit endpoint-policy denial.
    \item VPC flow logs (or S3 server access logs with \texttt{vpce} source) evidence the private path.
    \item A Route53/hosted check confirms no public DNS egress path exists from the subnet.
\end{itemize}

\section{Friday Use Case: A Healthcare Data Perimeter That Survives Credential Theft}
\index{use case!data perimeter}
A healthcare provider's threat model states: \emph{assume an IAM credential will eventually be compromised; no data may leave to an outside AWS account even then.} Design the perimeter for the analytics platform.

\subsection{The Layered Design}
\begin{figure}[H]
    \centering
    \resizebox{0.95\textwidth}{!}{%
    \begin{tikzpicture}[
        layer/.style={rectangle, rounded corners, draw=red!60!black, thick, fill=red!4, align=center, minimum width=3.6cm, minimum height=1.15cm, font=\footnotesize\bfseries},
        data/.style={cylinder, shape aspect=0.25, draw=primary, thick, fill=blue!6, shape border rotate=90, align=center, text width=2.6cm, minimum height=1.3cm, font=\small\bfseries},
        arr/.style={-{Stealth[length=2.5mm]}, draw=red!60!black, thick}
    ]
        \node[data] (lake) at (0,0) {Lake + warehouse\\(KMS-encrypted)};
        \node[layer] (id) at (5.0,2.2) {Identity perimeter\\roles only, no keys\\aws:PrincipalOrgID conditions};
        \node[layer] (net) at (5.0,0) {Network perimeter\\VPC endpoints only\\endpoint policies};
        \node[layer] (org) at (5.0,-2.2) {Organization perimeter\\SCPs: deny external share\\deny non-endpoint S3};
        \node[layer] (keys) at (-0.2,2.2) {Key perimeter\\KMS key policies\\org-only decrypt};
        \draw[arr] (id) -- (lake);
        \draw[arr] (net) -- (lake);
        \draw[arr] (org) -- (lake);
        \draw[arr] (keys) -- (lake);
    \end{tikzpicture}%
    }
    \caption{The four perimeters: identity, network, organization, and keys. A stolen credential must defeat all four to exfiltrate data.}
    \label{fig:ch21-perimeter}
\end{figure}

\textbf{Identity perimeter.} All access via roles with conditions requiring the organization (\texttt{aws:PrincipalOrgID}) and, where applicable, the expected network (\texttt{aws:SourceVpce}). No IAM users with long-lived keys exist to steal in the first place; workforce access is federated.

\textbf{Network perimeter.} Data services are reachable only through VPC endpoints whose policies allow only the organization's buckets and only expected actions. There is no internet gateway in data subnets. A credential used from an attacker's laptop satisfies no network condition.

\textbf{Organization perimeter (SCP).} An SCP denies \texttt{s3:PutObject} to any bucket where the destination account is outside the organization, denies disabling CloudTrail, and confines regions. Even a fully compromised account admin cannot exceed an SCP.

\textbf{Key perimeter.} Bucket-default KMS keys whose policies allow decrypt only for organization principals. Data copied out---if every other layer somehow failed---arrives as ciphertext without a usable key.

\begin{lstlisting}[language=json,caption={The anti-exfiltration SCP: deny S3 writes to non-organization accounts.},label={lst:ch21-scp}]
{
  "Effect": "Deny",
  "Action": ["s3:PutObject", "s3:PutObjectAcl"],
  "Resource": "arn:aws:s3:::*/*",
  "Condition": {
    "StringNotEquals": {
      "aws:PrincipalAccount": ["111111111111", "222222222222"]
    }
  }
}
\end{lstlisting}

\subsection{Verification and Operations}
The perimeter is tested like code: a red-team drill attempts exfiltration with a deliberately ``leaked'' sandbox credential and must fail at a named layer, with the CloudTrail event captured. Continuously, Config rules check for internet gateways in data VPCs and for buckets without org-scoped key policies, and CloudTrail Lake Formation/S3 data events alarm on access patterns outside the endpoint set.

\begin{pitfall}
\textbf{A perimeter that exists only in diagrams.} Each layer must have a test that would fail if the layer were removed. Untested SCPs have typos; untested endpoint policies have wildcards; untested key policies allow \texttt{*}. The red-team drill is the unit test of the perimeter.
\end{pitfall}

\section{Interview Q\&A}
\subsection{Concepts and Trade-offs}
\begin{enumerate}
    \item \textbf{Q: Gateway versus interface VPC endpoint: which is free for S3 traffic?}\\
    A: The gateway endpoint---a route-table entry for S3 and DynamoDB with no hourly or per-GB charge. Interface endpoints (PrivateLink) cover other services and bill per AZ-hour plus data.
    \item \textbf{Q: What does a permissions boundary limit on a role?}\\
    A: The role's maximum permissions regardless of its identity policies---the intersection of policy and boundary is effective. It makes delegated role creation safe.
    \item \textbf{Q: Why is automatic Secrets Manager rotation better than static keys?}\\
    A: Rotation bounds the value of any leaked credential to the rotation window and updates the target system automatically; static keys leak forever and rotate never.
    \item \textbf{Q: How do SCPs differ from IAM policies?}\\
    A: SCPs apply at the Organizations level and cap what any account---including its administrators---may do. IAM policies grant within an account; SCPs bound the account itself.
    \item \textbf{Q: What exfiltration path does PrivateLink close that security groups cannot?}\\
    A: Security groups control instance traffic but not which \emph{service endpoints} data flows to. Endpoint policies on PrivateLink/gateway endpoints restrict the destinations themselves---your network cannot be used to write to a stranger's bucket.
    \item \textbf{Q: Why does KMS complete the perimeter?}\\
    A: Key policies can restrict decryption to organization principals, so data copied past every other control arrives as unusable ciphertext.
\end{enumerate}

\section{Further Reading}
The IAM User Guide covers roles, boundaries, and policy evaluation \cite{awsIAMDocs}; the PrivateLink guide covers both endpoint types and endpoint policies \cite{awsPrivateLinkDocs}; KMS and Secrets Manager documentation cover envelope encryption, key policies, and rotation \cite{awsKMSDocs,awsSecretsManagerDocs}. The Config and CloudTrail guides support the verification layer \cite{awsConfigDocs,awsCloudTrailDocs}.

\section*{Key Takeaways}
\begin{itemize}
    \item Roles everywhere, users never; least privilege iterated from observed usage.
    \item Boundaries and SCPs are ceilings that make every future grant safer; deploy them first.
    \item Gateway endpoints make private S3 free; interface endpoints privatize everything else---with a price to watch.
    \item Endpoint policies are the network half of the perimeter, independent of IAM.
    \item KMS key policies are the final layer: ciphertext without an org key is not exfiltration.
    \item Perimeters are tested by drills, not diagrams; every layer needs a test that would fail without it.
\end{itemize}

\section{Exercises}
\begin{enumerate}
    \item \textbf{(Conceptual)} Walk the IAM evaluation logic for a role with an allow policy, a boundary, and an SCP. What wins in a conflict?
    \item \textbf{(Conceptual)} Why does the identity-perimeter condition on \texttt{aws:SourceVpce} defeat a stolen credential used from an attacker's machine?
    \item \textbf{(Hands-on)} Run the Thursday lab; then tighten the endpoint policy to read-only and demonstrate the write failing.
    \item \textbf{(Hands-on)} Create a Secrets Manager secret for a database, wire automatic rotation, and verify a rotation cycle.
    \item \textbf{(Hands-on)} Use IAM Access Analyzer to generate a least-privilege policy from a role's CloudTrail history and diff it against the hand-written policy.
    \item \textbf{(Design)} Write the SCP set for a three-account data estate: region restriction, CloudTrail protection, and exfiltration denial.
    \item \textbf{(Design)} Price the perimeter for the Friday scenario: endpoints per AZ, KMS request volume, Config rules, and CloudTrail data events.
    \item \textbf{(Architecture)} Design break-glass access for the perimeter: who can bypass, how is it alarmed, and how does it expire?
\end{enumerate}

\section{Capstone Project: A Tested Data Perimeter for the Analytics Estate}\label{sec:ch21-capstone}
\index{capstone project}\index{data perimeter!capstone}

\begin{capstonebox}
\textbf{Mission.} Wrap the course's analytics estate in a tested perimeter: org-scoped SCPs, endpoint-privatized networking, org-scoped KMS keys, rotated secrets---and prove it with a red-team drill that fails at a named layer.

\textbf{Estimated time.} 6--8 hours in a sandbox organization.

\textbf{What you\textquotesingle{}ll practice.}
\begin{itemize}
    \item SCP authoring and attachment with safe rollout.
    \item Gateway and interface endpoint design with endpoint policies.
    \item KMS key policy design for cross-account analytics.
    \item Evidence-driven verification: drills, Config rules, CloudTrail alarms.
\end{itemize}
\end{capstonebox}

\subsection*{Scenario}
The healthcare provider from Friday approves a proof-of-concept on two sandbox accounts (data and analytics). Your job: the estate must hold against a simulated credential leak, and every control must come with its test.

\subsection*{Requirements}
\textbf{Functional requirements.}
\begin{itemize}
    \item SCPs: region restriction, anti-exfiltration (\cref{lst:ch21-scp} generalized to \texttt{aws:PrincipalOrgID}), CloudTrail protection.
    \item S3 gateway endpoint plus interface endpoints for the services in use, with bucket-scoped endpoint policies.
    \item Bucket-default KMS with organization-scoped decrypt; Secrets Manager with rotation for the database credential.
    \item The red-team drill: a ``leaked'' sandbox credential attempting console and API exfiltration.
\end{itemize}

\textbf{Non-functional requirements.}
\begin{itemize}
    \item Every layer has a test that demonstrably fails when the control is removed (in the sandbox).
    \item Config rules and CloudTrail alarms watch the perimeter continuously.
    \item A cost statement covers the perimeter's own monthly bill.
\end{itemize}

\subsection*{Implementation Steps}
\begin{enumerate}
    \item Stand up the two-account organization; attach SCPs progressively and test each.
    \item Build the private networking per \cref{lst:ch21-endpoint}, extended to the used services.
    \item Re-key the buckets to org-scoped CMKs; migrate the database credential into Secrets Manager with rotation.
    \item Author the drill script and run it; capture the layer and CloudTrail event of each failure.
    \item Deploy the Config rules and alarms; document the cost statement.
\end{enumerate}

\subsection*{Deliverables}
\begin{itemize}
    \item SCPs, endpoint policies, and key policies as code.
    \item The drill report: attempted actions, blocking layer, evidence events.
    \item Config rule set and alarm definitions with one fired example.
    \item The perimeter cost statement and the operations runbook.
\end{itemize}

\subsection*{Acceptance Criteria}
\begin{itemize}
    \item The drill's exfiltration attempts all fail at named, evidenced layers.
    \item Cross-account analytics (the lake read from the analytics account) still works---the perimeter permits the intended path.
    \item Rotation demonstrably invalidates the previous database credential.
    \item Removing any single control in the sandbox is detected by its test.
\end{itemize}

\subsection*{Stretch Goals}
\begin{itemize}
    \item Add Macie (Chapter~23) to close the loop on what the perimeter protects.
    \item Implement network-perimeter exceptions as time-boxed, alarmed break-glass roles.
    \item Extend the drill to an insider scenario: an authorized analyst exceeding their normal pattern.
    \item Codify the whole perimeter in Terraform (Chapter~24) with policy tests in CI.
\end{itemize}
